\documentclass[12pt, a4paper]{letter}
\usepackage[utf8]{inputenc}
\usepackage[T1]{fontenc}
\usepackage{geometry}
\usepackage{amssymb}
\usepackage{hyperref}

\geometry{margin=1in}

\hypersetup{colorlinks=true, linkcolor=blue, urlcolor=blue}

\signature{Rafael D. De Paz \\ Sovereign Architect \\ \texttt{me@rdepaz.com}}
\address{Rafael D. De Paz \\ Sovereign Architect}

\begin{document}
\begin{letter}{Editorial Board}

\opening{Dear Editors,}

I am writing to submit the enclosed manuscript, entitled

\begin{center}
\textbf{``Isomorphic Cryptographic Tethering: Using Physical Objects as Thermodynamic Memory Locks''}
\end{center}

\noindent for consideration for publication.

\medskip

\textbf{Summary.} Advancing beyond traditional Hardware Security Modules (HSMs) and Physically Unclonable Functions (PUFs), this paper explores Isomorphic Cryptographic Tethering. We formalize how the literal thermodynamic and geometric topologies of specific physical objects (e.g., structured crystalline lattices or specialized jewelry) can be mapped identically to the hash functions of an access protocol. This creates a hard physical tether to reality, bounding digital information systems within strict thermodynamic constraints to prevent total systemic amnesia.

\medskip

This manuscript has not been previously published and is not under consideration
at any other journal. I confirm no conflicts of interest.

\textbf{Cryptographic Lineage \& Validation.} This mathematical substrate has been cryptographically sealed and tracked on the global Sovereign Master Ledger to prevent retroactive editing and to verify the source authorship. The immutable SHA-256 integrity checksums, formal PDF renderings, and lineage authorities can be verified at \url{https://rdepaz.com/research}.

\medskip

\closing{Respectfully submitted,}
\end{letter}
\end{document}
